\chapter{Radiación térmica: por qué brilla lo caliente}\label{cap:radiacion}

\objetivos{
\begin{itemize}
  \item Describir la radiación electromagnética con su longitud de onda y ubicar el infrarrojo
  medio y el térmico en el espectro.
  \item Usar la ley de Planck, la ley de Wien y la ley de Stefan--Boltzmann para comparar un
  fuego con el suelo que lo rodea.
  \item Explicar por qué los sensores de fuego miran alrededor de 4~\um{} y de 11~\um{}.
  \item Calcular una temperatura de brillo a partir de una radiancia y entender qué significa.
\end{itemize}}

\section{La luz como onda}

Todo lo que un satélite sabe de la Tierra le llega en forma de \termino{radiación
electromagnética}: energía que viaja por el espacio como una onda de campos eléctricos y
magnéticos. La luz visible es un tipo de radiación electromagnética; las ondas de radio, las
microondas, el infrarrojo y los rayos~X son otros. Todos son el mismo fenómeno y solo se
diferencian en una cosa: la \termino{longitud de onda}, que se escribe $\lambda$ y es la
distancia entre dos crestas consecutivas de la onda.

Toda radiación electromagnética viaja en el vacío a la misma velocidad,
$c \approx \SI{3.00e8}{\metre\per\second}$. Si las crestas están más juntas, pasan más crestas
por segundo. El número de crestas por segundo es la \termino{frecuencia} $\nu$, y las tres
cantidades están atadas por
\begin{equation}
  c = \lambda\,\nu .
\end{equation}

En teledetección las longitudes de onda se expresan casi siempre en \termino{micrómetros}
(\um), la millonésima parte de un metro. Un cabello humano mide unos 70~\um{} de grosor; la luz
visible va de 0,4~\um{} (violeta) a 0,7~\um{} (rojo). La figura~\ref{fig:espectro} muestra la
parte del espectro que nos interesa, con los nombres que usaremos todo el libro.

\begin{figure}[ht]
\centering
\begin{tikzpicture}[x=1cm,y=1cm]
  \def\lx#1{(ln(#1)-ln(0.3))/(ln(20)-ln(0.3))*14}
  \shade[left color=violet!70,right color=red!80,middle color=green!70] ({\lx{0.4}},0) rectangle ({\lx{0.7}},1);
  \fill[red!25] ({\lx{0.7}},0) rectangle ({\lx{3}},1);
  \fill[fuego!45] ({\lx{3}},0) rectangle ({\lx{5}},1);
  \fill[brasa!30] ({\lx{5}},0) rectangle ({\lx{8}},1);
  \fill[brasa!55] ({\lx{8}},0) rectangle ({\lx{14}},1);
  \fill[brasa!20] ({\lx{14}},0) rectangle ({\lx{20}},1);
  \fill[cenizaClaro] ({\lx{0.3}},0) rectangle ({\lx{0.4}},1);
  \draw[ceniza] (0,0) rectangle (14,1);
  \foreach \v in {0.4,0.7,1,2,3,4,5,8,11,14,20}{
    \draw ({\lx{\v}},0)--({\lx{\v}},-0.12) node[below,font=\footnotesize]{\v};
  }
  \node[font=\footnotesize] at (7,-0.8) {longitud de onda $\lambda$ (\um, escala logarítmica)};
  \node[font=\footnotesize,align=center] at ({\lx{0.53}},1.35) {visible};
  \node[font=\footnotesize,align=center] at ({\lx{1.45}},1.45) {infrarrojo cercano\\ y de onda corta};
  \node[font=\footnotesize\bfseries,align=center] at ({\lx{3.9}},1.45) {infrarrojo\\ medio};
  \node[font=\footnotesize\bfseries,align=center] at ({\lx{10.6}},1.45) {infrarrojo\\ térmico};
  \draw[very thick,black] ({\lx{3.96}},-1.35)--({\lx{3.96}},-1.05); \node[below,font=\footnotesize] at ({\lx{3.96}},-1.35) {canal de 4~\um};
  \draw[very thick,black] ({\lx{11}},-1.35)--({\lx{11}},-1.05); \node[below,font=\footnotesize] at ({\lx{11}},-1.35) {canal de 11~\um};
\end{tikzpicture}
\caption{La porción del espectro electromagnético que usa la detección de fuego. Las dos marcas
inferiores señalan los dos canales protagonistas de este libro.}
\label{fig:espectro}
\end{figure}

\section{Todo cuerpo caliente emite}

Cualquier objeto cuya temperatura esté por encima del cero absoluto emite radiación
electromagnética, simplemente por estar caliente. Los átomos y moléculas de la materia se agitan
más cuanto mayor es la temperatura, y esa agitación de cargas eléctricas genera ondas. A esto se
le llama \termino{radiación térmica}.

No la vemos en la vida diaria porque a temperaturas ordinarias casi toda cae en el infrarrojo,
que el ojo no detecta. Pero se nota: es el calor que sientes en la cara cerca de una fogata sin
tocarla. Si el objeto se calienta lo suficiente, parte de su emisión entra en el visible y
empieza a verse: la resistencia de una cocina eléctrica pasa de negra a roja, y un metal en la
fragua, de rojo a amarillo y a blanco.

Ese cambio de color contiene las dos ideas centrales del capítulo:

\begin{enumerate}
  \item al aumentar la temperatura, el objeto emite \textbf{mucha más} radiación en total;
  \item al aumentar la temperatura, la emisión se desplaza hacia \textbf{longitudes de onda más
  cortas}.
\end{enumerate}

Las temperaturas se medirán siempre en \termino{kelvin} (K), la escala absoluta: $T(\mathrm{K})
= T(^\circ\mathrm{C}) + 273{,}15$. Un suelo a 27~$^\circ$C está a 300~K. Las llamas de un
incendio de vegetación están típicamente entre 800 y 1200~K, y las brasas sin llama entre 600 y
850~K.

\section{Radiancia: la cantidad que mide un sensor}

Antes de escribir fórmulas hay que precisar \emph{qué} se mide. Un sensor mira hacia el suelo a
través de una pequeña abertura, dentro de un cono estrecho de direcciones, y solo deja pasar una
banda estrecha de longitudes de onda. La cantidad física que describe la energía que le llega en
esas condiciones es la radiancia espectral.

\begin{definicion}{radiancia espectral}
La radiancia espectral $L_\lambda$ es la potencia que emite una superficie por unidad de área,
por unidad de ángulo sólido (el <<tamaño>> del cono de direcciones) y por unidad de longitud de
onda. Su unidad habitual en teledetección es
\si{\watt\per\square\metre\per\steradian\per\micro\metre}.
\end{definicion}

La propiedad que la hace tan útil es que la radiancia de una superficie extensa \textbf{no
depende de la distancia} desde la que se la mira: al alejarse, el sensor recibe menos potencia
de cada punto, pero dentro de su cono entran proporcionalmente más puntos, y ambos efectos se
cancelan. Por eso una radiancia medida a 700~km de altura puede compararse directamente con la
que emite el suelo (salvo por lo que haga la atmósfera en el camino, tema del
capítulo~\ref{cap:atmosfera}).

\section{El cuerpo negro y la ley de Planck}

Para tener una referencia se define un emisor ideal.

\begin{definicion}{cuerpo negro}
Un cuerpo negro es un objeto ideal que absorbe toda la radiación que le llega y que, a cada
temperatura, emite la máxima radiación térmica posible. Su emisión depende únicamente de su
temperatura.
\end{definicion}

El nombre confunde: un cuerpo negro caliente brilla, y mucho. <<Negro>> solo indica que no
refleja nada. El Sol, una brasa y el suelo se comportan de forma bastante parecida a un cuerpo
negro.

La radiancia espectral de un cuerpo negro a temperatura $T$ la da la \termino{ley de Planck}:
\begin{equation}\label{eq:planck}
  B_\lambda(T) \;=\; \frac{2hc^2}{\lambda^5}\;\frac{1}{\exp\!\left(\dfrac{hc}{\lambda k T}\right)-1},
\end{equation}
donde $h = \SI{6.626e-34}{\joule\second}$ es la constante de Planck,
$k = \SI{1.381e-23}{\joule\per\kelvin}$ es la constante de Boltzmann y $c$ es la velocidad de la
luz. No hace falta memorizarla, pero sí saber leerla. Tiene dos factores. El primero,
$2hc^2/\lambda^5$, crece muy deprisa hacia longitudes de onda cortas. El segundo, el de la
exponencial, se desploma hacia longitudes de onda cortas cuando $\lambda T$ es pequeño. El
producto es una curva con forma de joroba: sube, alcanza un máximo y baja.

En la práctica se usa en una forma más cómoda, agrupando constantes:
\begin{equation}\label{eq:planck-practica}
  B_\lambda(T) = \frac{c_1}{\lambda^5\left[\exp\!\left(\dfrac{c_2}{\lambda T}\right)-1\right]},
  \qquad
  \begin{aligned}
  c_1 &= \SI{1.191e8}{\watt\per\square\metre\per\steradian\micro\metre\tothe{4}},\\
  c_2 &= \SI{1.4388e4}{\micro\metre\kelvin},
  \end{aligned}
\end{equation}
con $\lambda$ en micrómetros y $T$ en kelvin. El resultado sale directamente en
\si{\watt\per\square\metre\per\steradian\per\micro\metre}.

\begin{figure}[ht]
\centering
\begin{tikzpicture}
\begin{axis}[width=13.5cm,height=8cm,xmode=log,ymode=log,
  xmin=1,xmax=20,ymin=0.01,ymax=20000,
  xlabel={Longitud de onda $\lambda$ (\um)},
  ylabel={Radiancia $B_\lambda$ (\si{\watt\per\square\metre\per\steradian\per\micro\metre})},
  xtick={1,2,3,4,5,8,11,15,20},xticklabels={1,2,3,4,5,8,11,15,20},
  legend style={at={(0.5,-0.2)},anchor=north,font=\small,draw=none,legend columns=3,/tikz/every even column/.append style={column sep=8pt}},
  grid=major,grid style={ceniza!15},domain=1:20,samples=160,
  declare function={planck(\l,\T)=1.191e8/(\l^5*(exp(14388/(\l*\T))-1));}]
\addplot[fuego,very thick] {planck(x,1200)};
\addplot[brasa,very thick] {planck(x,1000)};
\addplot[brasa!60!black,very thick,dashed] {planck(x,800)};
\addplot[violeta,very thick,dashdotted] {planck(x,600)};
\addplot[cielo,very thick] {planck(x,300)};
\legend{1200 K (llama intensa),1000 K (llama),800 K,600 K (brasa),300 K (suelo)}
\draw[black,thick,dashed] (axis cs:3.96,0.01)--(axis cs:3.96,20000);
\draw[black,thick,dashed] (axis cs:11,0.01)--(axis cs:11,20000);
\node[font=\footnotesize,rotate=90,anchor=south east] at (axis cs:3.96,15000) {4~\um};
\node[font=\footnotesize,rotate=90,anchor=south east] at (axis cs:11,15000) {11~\um};
\draw[<->,thick] (axis cs:4.25,0.6726)--(axis cs:4.25,3320) node[midway,right,font=\footnotesize,fill=white,inner sep=1pt]{$\times 4900$};
\draw[<->,thick] (axis cs:12,9.573)--(axis cs:12,274) node[midway,right,font=\footnotesize,fill=white,inner sep=1pt]{$\times 29$};
\end{axis}
\end{tikzpicture}
\caption{Curvas de Planck para el suelo (300~K) y para varias temperaturas de fuego, en escala
logarítmica. A 4~\um{} una superficie a 1000~K emite unas 4900 veces más que una a 300~K; a
11~\um, solo unas 29 veces más. Esa asimetría es el fundamento de la detección de fuego.}
\label{fig:planck}
\end{figure}

\begin{ejemplo}{radiancia del suelo y del fuego a 4 micrómetros}
Calcula la radiancia de un cuerpo negro a $\lambda = \SI{3.96}{\micro\metre}$ para el suelo a
300~K y para una llama a 1000~K.
\solucion Con la ecuación~\eqref{eq:planck-practica}, $\lambda^5 = 3{,}96^5 = 973{,}8$.

Para 300~K: $c_2/(\lambda T) = 14388/(3{,}96 \times 300) = 12{,}11$, y $e^{12{,}11} - 1 =
\num{1.818e5}$. Entonces
\[
  B = \frac{\num{1.191e8}}{973{,}8 \times \num{1.818e5}} = \SI{0.673}{\watt\per\square\metre\per\steradian\per\micro\metre}.
\]
Para 1000~K: $c_2/(\lambda T) = 3{,}633$, y $e^{3{,}633} - 1 = 36{,}84$. Entonces
\[
  B = \frac{\num{1.191e8}}{973{,}8 \times 36{,}84} = \SI{3320}{\watt\per\square\metre\per\steradian\per\micro\metre}.
\]
La llama emite $3320/0{,}673 \approx 4900$ veces más que el suelo en esta longitud de onda.
\end{ejemplo}

\section{La ley de Wien: dónde está el máximo}

La joroba de Planck tiene su máximo en una longitud de onda que depende solo de la temperatura.
Derivando la ecuación~\eqref{eq:planck} e igualando a cero se obtiene la \termino{ley de
desplazamiento de Wien}:
\begin{equation}\label{eq:wien}
  \lambda_{\text{máx}} = \frac{\SI{2898}{\micro\metre\kelvin}}{T}.
\end{equation}
Cuanto más caliente es el cuerpo, más corta es la longitud de onda en la que más emite. La
tabla~\ref{tab:wien} aplica la ley a los objetos de este libro.

\begin{table}[ht]
\centering
\begin{tabular}{lrrl}
\toprule
Objeto & $T$ (K) & $\lambda_{\text{máx}}$ (\um) & Región del espectro \\
\midrule
Sol (superficie) & 5800 & 0,50 & visible \\
Llama intensa & 1200 & 2,4 & infrarrojo de onda corta \\
Llama & 1000 & 2,9 & infrarrojo de onda corta/medio \\
Brasa & 600 & 4,8 & infrarrojo medio \\
Suelo & 300 & 9,7 & infrarrojo térmico \\
Cima de nube alta & 220 & 13,2 & infrarrojo térmico \\
\bottomrule
\end{tabular}
\caption{Longitud de onda de máxima emisión según la ley de Wien.}
\label{tab:wien}
\end{table}

La tabla explica la elección de los dos canales. El suelo y las nubes emiten sobre todo en torno
a 10~\um: el canal de 11~\um{} es la ventana natural para medir la temperatura del fondo. El
fuego emite sobre todo entre 2,5 y 5~\um: el canal de 4~\um{} está en plena zona de emisión del
fuego y, a la vez, en una zona donde el suelo emite muy poco.

\section{La ley de Stefan--Boltzmann: cuánto en total}

Si se suma la emisión de un cuerpo negro en todas las longitudes de onda y en todas las
direcciones, se obtiene la potencia total que emite cada metro cuadrado de su superficie. El
resultado es la \termino{ley de Stefan--Boltzmann}:
\begin{equation}\label{eq:sb}
  M = \sigma T^4, \qquad \sigma = \SI{5.67e-8}{\watt\per\square\metre\per\kelvin\tothe{4}}.
\end{equation}
La dependencia con la cuarta potencia es brutal. Un metro cuadrado de suelo a 300~K emite
$\sigma \cdot 300^4 \approx \SI{459}{\watt}$. Un metro cuadrado de llama a 1000~K emite
$\sigma \cdot 1000^4 \approx \SI{56700}{\watt}$, es decir, $(1000/300)^4 \approx 123$ veces más.

\begin{intuicion}[Total frente a canal]
La ley de Stefan--Boltzmann dice que el fuego emite unas 120 veces más que el suelo
\emph{sumando todo el espectro}. La figura~\ref{fig:planck} dice que a 4~\um{} emite unas 4900
veces más. No hay contradicción: al calentarse, el cuerpo no solo emite más, sino que traslada
su emisión hacia las ondas cortas. Un canal bien elegido aprovecha los dos efectos a la vez, y
por eso el contraste en un canal puede ser mucho mayor que el contraste total.
\end{intuicion}

\section{Emisividad: los cuerpos reales}

Ningún objeto real es un cuerpo negro perfecto. La relación entre lo que emite un cuerpo real y
lo que emitiría un cuerpo negro a la misma temperatura se llama \termino{emisividad},
$\varepsilon$, y vale entre 0 y 1:
\begin{equation}
  L_\lambda = \varepsilon_\lambda \, B_\lambda(T).
\end{equation}
En el infrarrojo térmico la vegetación, el agua y la mayoría de los suelos tienen emisividades
entre 0,95 y 0,99, de modo que tratarlos como cuerpos negros es una buena aproximación. En el
infrarrojo medio la emisividad de los suelos secos y desnudos puede ser bastante menor, lo que
significa, como se verá en el capítulo~\ref{cap:atmosfera}, que \emph{reflejan} más luz solar
en ese canal. Esa es una de las fuentes de falsas alarmas.

\section{Temperatura de brillo}\label{sec:tb}

Los datos de focos no publican radiancias, sino temperaturas. Las columnas se llaman
\texttt{brightness}, \texttt{bright\_t31}, \texttt{bright\_ti4}, \texttt{bright\_ti5}, y sus
valores están en kelvin. Pero no son la temperatura de las llamas.

\begin{definicion}{temperatura de brillo}
La temperatura de brillo $\Tb$ de un píxel en un canal es la temperatura que tendría un cuerpo
negro que emitiera exactamente la radiancia medida en ese canal. Se obtiene invirtiendo la ley
de Planck:
\begin{equation}\label{eq:tb}
  \Tb = \frac{c_2}{\lambda\,\ln\!\left(1 + \dfrac{c_1}{\lambda^5 L_\lambda}\right)} .
\end{equation}
\end{definicion}

La temperatura de brillo es, en el fondo, otra forma de escribir la radiancia: una escala más
intuitiva. Si el píxel es una superficie uniforme con emisividad cercana a 1 y no hay atmósfera
de por medio, $\Tb$ coincide con su temperatura real. Si el píxel es una mezcla de suelo frío y
un poco de fuego muy caliente, $\Tb$ es un valor intermedio que no corresponde a la temperatura
de ninguna de las dos cosas.

\begin{ejemplo}{de radiancia a temperatura de brillo}
Un píxel muestra a $\lambda = \SI{3.96}{\micro\metre}$ una radiancia de
\SI{1.00}{\watt\per\square\metre\per\steradian\per\micro\metre}. ¿Cuál es su temperatura de
brillo?
\solucion Con la ecuación~\eqref{eq:tb}: $c_1/(\lambda^5 L) = \num{1.191e8}/(973{,}8 \times
1{,}00) = \num{1.223e5}$. Su logaritmo natural es $\ln(\num{1.223e5} + 1) = 11{,}71$. Entonces
\[
  \Tb = \frac{14388}{3{,}96 \times 11{,}71} = \SI{310.2}{\kelvin}.
\]
Comparando con el ejemplo anterior: pasar de 0,673 a 1,00 unidades de radiancia (un aumento del
49\,{}\%) equivale a subir de 300~K a 310~K. A 4~\um{} un cambio pequeño de temperatura produce un
cambio grande de radiancia.
\end{ejemplo}

\section{La sensibilidad a la temperatura de cada canal}

La última observación del ejemplo merece una sección, porque es la razón profunda por la que se
usan dos canales. Cerca de una temperatura dada, la radiancia de Planck se comporta
aproximadamente como una potencia de la temperatura, $B_\lambda \propto T^{\,n}$, con un
exponente
\begin{equation}
  n \approx \frac{c_2}{\lambda T}
\end{equation}
(válido cuando $c_2/\lambda T$ es bastante mayor que 1). Para un fondo a 300~K:
\[
  n_{4\,\mu\text{m}} \approx \frac{14388}{3{,}96\times 300} \approx 12,
  \qquad
  n_{11\,\mu\text{m}} \approx \frac{14388}{11\times 300} \approx 4{,}4 .
\]
A 4~\um{} la radiancia crece como la temperatura elevada a la potencia doce; a 11~\um, como la
temperatura a la potencia cuatro y pico. El canal de 4~\um{} es, por decirlo así, hipersensible
a cualquier cosa caliente. El de 11~\um{} es mucho más tranquilo. Un fuego pequeño dentro de un
píxel dispara el primero y apenas mueve el segundo. La \emph{diferencia} entre las temperaturas
de brillo de los dos canales, $\Delta T = T_4 - T_{11}$, es por tanto una firma del fuego, y es
la variable central del algoritmo de detección (capítulo~\ref{cap:deteccion}).

\begin{figure}[ht]
\centering
\begin{tikzpicture}
\begin{axis}[width=12.5cm,height=6.6cm,ymode=log,xmin=280,xmax=1200,ymin=0.5,ymax=20000,
  xlabel={Temperatura del cuerpo negro (K)},
  ylabel={Radiancia relativa a la de 300 K},
  legend style={at={(0.03,0.96)},anchor=north west,font=\small,draw=ceniza!40},
  grid=major,grid style={ceniza!15},domain=280:1200,samples=120,
  declare function={planck(\l,\T)=1.191e8/(\l^5*(exp(14388/(\l*\T))-1));}]
\addplot[fuego,very thick] {planck(3.96,x)/planck(3.96,300)};
\addplot[cielo,very thick] {planck(11,x)/planck(11,300)};
\legend{canal de 4~\um,canal de 11~\um}
\end{axis}
\end{tikzpicture}
\caption{Cuánto aumenta la radiancia de cada canal al calentar un cuerpo negro desde 300~K. El
canal de 4~\um{} se dispara; el de 11~\um{} crece con moderación.}
\label{fig:sensibilidad}
\end{figure}

\begin{parareto}
Las columnas de temperatura de brillo de los datos de FIRMS (capítulo~\ref{cap:firms}) son
exactamente las $T_4$ y $T_{11}$ de este capítulo: \texttt{brightness} y \texttt{bright\_t31}
en MODIS, \texttt{bright\_ti4} y \texttt{bright\_ti5} en VIIRS. No representan la temperatura de
las llamas y no son comparables entre sensores, porque dependen del tamaño del píxel. No las
uses como medida de <<intensidad>> común a MODIS y VIIRS.
\end{parareto}

\resumen

\begin{itemize}
  \item Todo cuerpo emite radiación térmica; cuanto más caliente, más emite y a longitudes de
  onda más cortas.
  \item La \textbf{radiancia espectral} es la cantidad que mide un sensor. No depende de la
  distancia a la superficie.
  \item La \textbf{ley de Planck} da la radiancia de un cuerpo negro; la de \textbf{Wien}, la
  longitud de onda del máximo ($\lambda_{\text{máx}} = 2898/T$); la de
  \textbf{Stefan--Boltzmann}, la potencia total ($\sigma T^4$).
  \item El suelo (300~K) emite sobre todo cerca de 10~\um; el fuego (600--1200~K), entre 2,5 y
  5~\um.
  \item A 4~\um{} una llama a 1000~K emite unas 4900 veces más que el suelo; a 11~\um, unas 29.
  \item La \textbf{temperatura de brillo} es la radiancia expresada como temperatura de cuerpo
  negro equivalente. En un píxel con mezcla de fuego y fondo no es la temperatura de ninguno.
  \item La diferencia $\Delta T = T_4 - T_{11}$ es la firma básica del fuego.
\end{itemize}

\ejercicios

\begin{enumerate}
  \item Usa la ley de Wien para hallar la longitud de onda de máxima emisión de una brasa a
  700~K y de un cuerpo humano a 310~K.
  \item Calcula con la ecuación~\eqref{eq:planck-practica} la radiancia a 11~\um{} de un cuerpo
  negro a 300~K y a 1000~K, y comprueba que el cociente es cercano a 29.
  \item ¿Cuántas veces más potencia total emite por metro cuadrado una llama a 1200~K que una
  brasa a 600~K?
  \item Un píxel tiene a 11~\um{} una radiancia de
  \SI{12.0}{\watt\per\square\metre\per\steradian\per\micro\metre}. Calcula su temperatura de
  brillo.
  \item Explica por qué un canal en el visible (0,65~\um) no sirve para detectar fuego de día
  aunque las llamas se vean a simple vista. Pista: compara, con la ley de Planck, lo que emite
  una llama a 1000~K en esa longitud de onda con lo que sugiere la tabla~\ref{tab:wien} sobre
  el Sol.
  \item Repite el cálculo del exponente $n \approx c_2/\lambda T$ para un fondo a 320~K (un
  suelo desnudo al mediodía). ¿Aumenta o disminuye la ventaja del canal de 4~\um{} en un día
  caluroso?
\end{enumerate}

\lecturas

El capítulo sobre radiación térmica de cualquier texto de física universitaria (por ejemplo, la
sección de radiación de cuerpo negro de Sears y Zemansky) desarrolla estas leyes con más
detalle. Para su uso en teledetección, el manual de usuario del producto de fuego de MODIS
(apéndice~\ref{ap:fuentes}) resume las bandas y sus temperaturas de saturación.
